\documentclass[12pt,a4paper]{article}
\usepackage[utf8]{inputenc}
\usepackage[T1]{fontenc}
\usepackage{amsmath, amsfonts, amssymb}
\usepackage{geometry}
\usepackage{lmodern}
\usepackage{tcolorbox}

\geometry{hmargin=2cm,vmargin=1.5cm}
\setlength{\parindent}{0pt}
\setlength{\parskip}{5pt}

\begin{document}

\begin{center}
    \rule{\linewidth}{0.5pt} \\
    \Large \textbf{Equation de Korteweg-de Vries} \\
    \rule{\linewidth}{0.5pt} \\
    \vspace{0.5cm}
\end{center}

On a notre EDP :
\begin{equation}
  \frac{\partial u(x,t)}{\partial t} 
  + \epsilon u(x,t)\frac{\partial u(x,t)}{\partial x}
  + \mu \frac{\partial^{3}u(x,t)}{\partial x^{3}} = 0
  \label{eq:initiale}
\end{equation}

\vspace{0.5cm}
On pose $\xi=x-ct$
\begin{equation}
  u(x,t)=z(x-ct)=z(\xi)
  \label{eq:z}
\end{equation}

\vspace{0.5cm}
Calcul des dérivées :
\begin{align*}
\frac{\partial u}{\partial t} =\frac{dz}{dt}
&=\frac{dz}{d(x-ct)}\frac{\partial(x-ct)}{\partial t}=-c\frac{dz}{d\xi} \\
\frac{\partial u}{\partial x}=\frac{dz}{dx}
&=\frac{dz}{d(x-ct)}\frac{\partial(x-ct)}{\partial x}=\frac{dz}{d\xi} \\
\frac{\partial^{3}u}{\partial x^{3}} &=
\left(\frac{\partial}{\partial x}\right)^{3}u=\frac{d^{3}z}{d\xi^{3}}
\end{align*}

On substitue (\ref{eq:z}) à (\ref{eq:initiale}) et on obtient :
\[ \boxed{-c\frac{dz(\xi)}{d\xi}+\epsilon z(\xi)\frac{dz(\xi)}{d\xi}+\mu\frac{d^{3}z(\xi)}{d\xi^{3}}=0} \]

Transformation de la dérivée totale en dérivée partielle (intégration) :
\begin{align*}
  C_{1} &=-c\int dz+\epsilon\int z dz+\mu\int\frac{d^{3}z(\xi)}{d^{3}\xi} d\xi \\ 
  C_{1} &=-cz(\xi)+\frac{\epsilon}{2}z^{2}(\xi)+\mu\frac{d^{2}z(\xi)}{d\xi^{2}}
\end{align*}

On multiplie par $\frac{dz}{d\xi}$ :
\begin{align*}
  C_{1}\frac{dz}{d\xi} &=-cz\frac{dz}{d\xi}+\frac{\epsilon}{2}z^{2}\frac{dz}{d\xi}
  +\mu\frac{d^{2}z}{d\xi^{2}}\cdot\frac{dz}{d\xi} \\
  C_{1}dz &=-cz dz+\frac{\epsilon}{2}z^{2}dz+\mu\frac{d^{2}z}{d\xi^{2}}dz
\end{align*}

\newpage{}

On intègre des deux côtés et on pose une constante d'intégration $C_2$ :
\begin{align*}
  C_{1}\int dz &=-c\int z dz+\frac{\epsilon}{2}\int z^{2}dz
  +\mu\int\frac{d^{2}z}{d\xi^{2}}dz \\
  C_{1}z+C_{2} &=-\frac{c}{2}z^{2}+\frac{\epsilon}{6}z^{3}
  +\frac{1}{2}\mu\left(\frac{dz}{d\xi}\right)^{2}
\end{align*}

Lorsque $x \to \pm \infty$ on a : $z \to 0$ ; $\frac{dz}{d\xi} \to 0$ ; $\frac{d^{2}z}{d\xi^{2}}\to 0$

Ce qui implique que $C_{1}=C_{2}=0$.


Ainsi :
\begin{align}
  0 &=-\frac{c}{2}z^{2}+\frac{\epsilon}{6}z^{3}+\frac{1}{2}\mu\left(\frac{dz}{d\xi}\right)^{2} \notag \\
  -\frac{1}{2}\mu\left(\frac{dz}{d\xi}\right)^{2} &=-\frac{c}{2}z^{2}+\frac{\epsilon}{6}z^{3} \notag \\
  \left(\frac{dz}{d\xi}\right)^{2} 
  &=\frac{c}{\mu}z^{2}-\frac{\epsilon}{3\mu}z^{3} \notag\\
  \left(\frac{dz}{d\xi}\right)^{2} 
  &=z^{2}\left(\frac{c}{\mu}-\frac{\epsilon}{3\mu}z\right)\label{eq:dz2}
\end{align}

On a donc :
\[ \boxed{\frac{dz}{d\xi} = z\sqrt{\frac{c}{\mu}-\frac{\epsilon}{3\mu}z}} \]

On sépare les variables $z$ et $\xi$ :
\begin{equation}
  \frac{dz}{z\sqrt{\frac{c}{\mu}-\frac{\epsilon}{3\mu}z}} = d\xi
  \label{eq:separee}
\end{equation}

On effectue un changement de variable pour simplifier le terme sous la racine. 

On pose :
\begin{equation}
  \begin{cases}
    z &= \frac{3\mu}{\epsilon}g^2 \\
    dz &= \frac{6\mu}{\epsilon}g \, dg 
  \end{cases}
  \label{eq:intro_g}
\end{equation}

\begin{tcolorbox}[
  colback=black!5!white,
  colframe=white!50!black,
  title=Rappel,
  fonttitle=\bfseries
]
\begin{equation}
  \int \frac{dx}{x\sqrt{a^{2} - x^{2}}} = \frac{1}{a} sech^{-1} \Big(\frac{x}{a}\Big)
  \label{eq:rappel}
\end{equation}
\end{tcolorbox}


On substitue $z$ et $dz$ dans \eqref{eq:separee} :
\begin{align}
\frac{\frac{6\mu}{\epsilon}g \, dg}{\frac{3\mu}{\epsilon}g^2 \sqrt{\frac{c}{\mu}
-\frac{\epsilon}{3\mu}\left(\frac{3\mu}{\epsilon}g^2\right)}} &= d\xi \notag\\
\frac{2 \, dg}{g \sqrt{\frac{c}{\mu}-g^2}} &= d\xi \label{eq:eq_a_int}
\end{align}

\newpage

On note $\xi_{0}$ le point auqel $z(\xi_{0})=\max{z}$. Ainsi,
d'après \eqref{eq:intro_g}, $g(\xi_{0})=\max{g}$\\

Calculons ce maximum. On utilise \eqref{eq:dz2} et le maximum est atteint 
lorsque $\frac{dz}{d\xi}=0$. On suppose $z\ne 0$.

\begin{align*}
  z_{max}^{2}(\frac{c}{\mu}-\frac{\epsilon}{3\mu}z_{max}) &=0 \\
  \frac{c}{\mu}-\frac{\epsilon}{3\mu}z_{max} &=0 \\
  z_{max} &= \frac{3c}{\epsilon}
\end{align*}

On cherche la valeur correspondante pour $g$

\begin{align*}
  z_{max} &=\frac{3\mu}{\epsilon}g_{max}^{2} \\
  z_{max} \frac{\epsilon}{3\mu} &= g_{max} \\
  \frac{c}{\mu} &= g_{max}^{2} \\
  \sqrt{\frac{c}{\mu}} &= g_{max}
\end{align*}

On intègre \eqref{eq:eq_a_int} des deux côtés :
$$
2 \int_{\sqrt{\frac{c}{\mu}}}^{g(\xi)} 
\frac{d\eta}{\eta \sqrt{\frac{c}{\mu}-\eta^2}} = \int_{\xi_{0}}^{\xi} d\zeta
$$
On utilise \eqref{eq:rappel}
$$
2 \left[ - \frac{1}{\sqrt{\frac{c}{\mu}}} 
\text{sech}^{-1}\left(\frac{\eta}{\sqrt{\frac{c}{\mu}}}\right) \right]^{\eta=g(\xi)}_{\eta=\sqrt{\frac{c}{\mu}}} = (\xi - \xi_0)
$$

Sachant que $\text{sech}^{-1}=0$, on a :

\begin{align}
  2 \left( - \frac{1}{\sqrt{\frac{c}{\mu}}}
  \text{sech}^{-1}\left(\frac{g(\xi)}{\sqrt{\frac{c}{\mu}}}\right) \right) &= (\xi - \xi_0) \notag\\ 
  \text{sech}^{-1}\left(\frac{g(\xi)}{\sqrt{\frac{c}{\mu}}}\right) &= -\frac{(\xi - \xi_0)}{2} \sqrt{\frac{c}{\mu}} \notag\\
  \frac{g(\xi)}{\sqrt{\frac{c}{\mu}}} &= \text{sech}\left( \frac{(\xi - \xi_0)}{2} \sqrt{\frac{c}{\mu}} \right) \notag\\
g(\xi) &=\boxed{\sqrt{\frac{c}{\mu}} \text{sech}\left( \frac{(\xi - \xi_0)}{2} \sqrt{\frac{c}{\mu}} \right)} \label{eq:g_final}
\end{align}

\newpage{}

On substitue \eqref{eq:g_final} à \eqref{eq:intro_g} et on obtient :
\[ z(\xi) = \frac{3\mu}{\epsilon} \left( \sqrt{\frac{c}{\mu}} \text{sech}\left( \frac{(\xi - \xi_0)}{2} \sqrt{\frac{c}{\mu}} \right) \right)^2 \]
La solution finale est donc :
\[ \boxed{z(\xi) = \frac{3c}{\epsilon} \text{sech}^2\left( \frac{1}{2}\sqrt{\frac{c}{\mu}}(\xi - \xi_0) \right)} \]

\end{document}
